\documentclass[10pt,letterpaper]{article}

\usepackage[margin=1.6cm]{geometry}
\usepackage{fontspec}
\usepackage{xcolor}
\usepackage{amssymb}
\usepackage{tabularx}
\usepackage{array}
\usepackage{enumitem}
\usepackage{titlesec}
\usepackage{fancyhdr}
\usepackage{lastpage}
\usepackage{ragged2e}
\usepackage{listings}
\usepackage[most]{tcolorbox}
\usepackage[hidelinks]{hyperref}

\setmainfont{Arial}
\setsansfont{Arial}
\setmonofont{Consolas}

\definecolor{AIEPRed}{HTML}{D62027}
\definecolor{AIEPNavy}{HTML}{102A43}
\definecolor{AIEPInk}{HTML}{243B53}
\definecolor{AIEPSlate}{HTML}{52606D}
\definecolor{AIEPPaper}{HTML}{F8F3EC}
\definecolor{AIEPSoftBlue}{HTML}{E6EEF7}
\definecolor{AIEPGreen}{HTML}{3FAE6A}
\definecolor{AIEPGold}{HTML}{E0BC5A}
\definecolor{AIEPAmber}{HTML}{E8A33D}
\definecolor{Line}{HTML}{D8E1EA}

\pagestyle{fancy}
\fancyhf{}
\renewcommand{\headrulewidth}{0pt}
\fancyfoot[L]{\small\textcolor{AIEPSlate}{GeoGreen · Guía de armado del prototipo}}
\fancyfoot[R]{\small\textcolor{AIEPSlate}{\thepage/\pageref{LastPage}}}

\setlength{\parindent}{0pt}
\setlength{\parskip}{4pt}
\setlist[itemize]{leftmargin=*,topsep=1pt,itemsep=1pt}
\setlist[enumerate]{leftmargin=*,topsep=1pt,itemsep=1pt}
\renewcommand{\arraystretch}{1.3}

\titleformat{\section}{\large\bfseries\color{AIEPNavy}}{}{0pt}{}
\titlespacing{\section}{0pt}{12pt}{4pt}
\titleformat{\subsection}{\bfseries\color{AIEPInk}}{}{0pt}{}
\titlespacing{\subsection}{0pt}{8pt}{2pt}

\newcolumntype{Y}{>{\RaggedRight\arraybackslash}X}
\newcolumntype{L}[1]{>{\RaggedRight\arraybackslash}p{#1}}

\newtcolorbox{nota}{enhanced, colback=AIEPSoftBlue, colframe=AIEPNavy,
  boxrule=0pt, leftrule=2.5pt, arc=0.8mm, left=3mm, right=3mm, top=2mm, bottom=2mm,
  before skip=4pt, after skip=4pt}
\newtcolorbox{ojo}{enhanced, colback=AIEPRed!6, colframe=AIEPRed,
  boxrule=0pt, leftrule=2.5pt, arc=0.8mm, left=3mm, right=3mm, top=2mm, bottom=2mm,
  before skip=4pt, after skip=4pt}

\lstdefinestyle{plano}{
  basicstyle=\ttfamily\scriptsize\color{AIEPInk},
  frame=single, rulecolor=\color{Line}, backgroundcolor=\color{AIEPPaper},
  xleftmargin=3mm, xrightmargin=2mm, aboveskip=5pt, belowskip=5pt
}
\lstdefinestyle{cmd}{
  basicstyle=\ttfamily\footnotesize, frame=single, rulecolor=\color{Line},
  backgroundcolor=\color{AIEPSoftBlue!35},
  xleftmargin=3mm, xrightmargin=2mm, aboveskip=5pt, belowskip=5pt
}

% Chip de pin (recuadro oscuro tipo serigrafia)
\newcommand{\pin}[1]{\tcbox[on line, colback=AIEPInk, colframe=AIEPInk, boxrule=0pt,
  arc=0.7mm, left=1.4mm, right=1.4mm, top=0.3mm, bottom=0.3mm]%
  {\footnotesize\bfseries\textcolor{white}{\texttt{#1}}}}

\begin{document}

% ===================== PORTADA =====================
\begin{tcolorbox}[enhanced, colback=AIEPNavy, colframe=AIEPNavy, arc=0mm,
  boxrule=0pt, left=5mm, right=5mm, top=5mm, bottom=5mm]
{\color{white}\Huge\bfseries Guía de armado del prototipo}\\[2mm]
{\color{white}\Large GeoGreen · Monitor de llenado con Arduino}\\[3mm]
{\color{AIEPSoftBlue} Mapa de pines, conexiones y firmware para volver a armar o replicar el
prototipo desde cero.}
\end{tcolorbox}

\begin{nota}
\textbf{Qué es.} Un monitor de llenado de contenedor: un sensor mide la distancia al contenido,
el firmware la convierte en \textbf{porcentaje de llenado}, un \textbf{semáforo} de LEDs lo indica
y un \textbf{buzzer} alerta cuando está lleno. Una \textbf{pantalla OLED} muestra el porcentaje y
la \textbf{matriz LED} de la placa muestra la marca del proyecto.
Lógica: \textbf{medir $\rightarrow$ indicar $\rightarrow$ alertar}.
\end{nota}

% ===================== COMPONENTES =====================
\section{Componentes}
\begin{tabularx}{\linewidth}{@{}L{5.2cm}Y@{}}
\hline
\textbf{Componente} & \textbf{Especificación}\\
\hline
Placa controladora & Arduino UNO R4 WiFi (5\,V, matriz LED 12$\times$8 integrada)\\
Protoboard & 830 puntos (MB-102)\\
Sensor & HC-SR04 (ultrasónico, 4 pines)\\
LED 5\,mm $\times$3 & Verde, amarillo y rojo\\
Resistencias $\times$3 & 220\,$\Omega$ (o 330\,$\Omega$), una por LED\\
Buzzer & Activo 5\,V (2 pines, con polaridad)\\
Pantalla & OLED 0.96" I2C, SSD1306 (4 pines)\\
Cables & Jumper dupont macho-macho\\
Alimentación & Cable USB (desde el PC o un cargador 5\,V)\\
\hline
\end{tabularx}

\begin{nota}
En una \textbf{Arduino UNO R3} también funciona (mismo cableado y firmware); solo no tendría la
matriz LED integrada para la marca.
\end{nota}

% ===================== MAPA DE PINES =====================
\section{Mapa de pines (conexiones)}
La placa se alimenta por \textbf{USB}. Lleva el \pin{5V} de la placa al \textbf{riel rojo (+)} y
el \pin{GND} al \textbf{riel azul ($-$)} de la protoboard; de ahí salen las alimentaciones.

\begin{tabularx}{\linewidth}{@{}L{3.6cm}L{3.6cm}Y@{}}
\hline
\textbf{Componente} & \textbf{Pin del componente} & \textbf{Va a}\\
\hline
Sensor HC-SR04 & VCC & riel rojo (+)\\
 & GND & riel azul ($-$)\\
 & Trig & \pin{D9}\\
 & Echo & \pin{D10} \quad(directo, sin divisor)\\
\hline
LED rojo & ánodo (+, pata larga) & \pin{D2} a través de resistencia\\
 & cátodo ($-$, pata corta) & riel azul ($-$)\\
\hline
LED amarillo & ánodo (+) & \pin{D3} a través de resistencia\\
 & cátodo ($-$) & riel azul ($-$)\\
\hline
LED verde & ánodo (+) & \pin{D4} a través de resistencia\\
 & cátodo ($-$) & riel azul ($-$)\\
\hline
Buzzer & + (pata larga) & \pin{D8}\\
 & $-$ (pata corta) & riel azul ($-$)\\
\hline
Pantalla OLED & VCC & riel rojo (+)\\
 & GND & riel azul ($-$)\\
 & SDA & \pin{SDA} \quad(o \texttt{A4})\\
 & SCL & \pin{SCL} \quad(o \texttt{A5})\\
\hline
Matriz LED 12$\times$8 & --- & integrada en la placa (sin cableado)\\
\hline
\end{tabularx}

\subsection{Resumen rápido del Arduino}
\begin{tabularx}{\linewidth}{@{}L{2.6cm}Y@{}}
\hline
\pin{D2} & LED rojo (vía resistencia) \\
\pin{D3} & LED amarillo (vía resistencia) \\
\pin{D4} & LED verde (vía resistencia) \\
\pin{D8} & Buzzer (+) \\
\pin{D9} & HC-SR04 Trig \\
\pin{D10} & HC-SR04 Echo \\
\pin{SDA} / \pin{SCL} & Pantalla OLED (I2C) \\
\pin{5V} / \pin{GND} & Rieles + y $-$ de la protoboard \\
\hline
\end{tabularx}

% ===================== ESQUEMA =====================
\section{Esquema de conexión}
\begin{lstlisting}[style=plano]
   ARDUINO UNO R4 WiFi   (USB al PC = alimentacion 5V)
   |
   5V  ---> riel rojo (+)        riel + alimenta sensor y OLED
   GND ---> riel azul (-)        riel - es la tierra comun
   |
   D9  ---> HC-SR04 Trig
   D10 ---> HC-SR04 Echo         (directo, la placa es 5V)
   D2  --[R]--> LED rojo (+)     ; (-) -> riel (-)
   D3  --[R]--> LED amarillo (+) ; (-) -> riel (-)
   D4  --[R]--> LED verde (+)    ; (-) -> riel (-)
   D8  -------> Buzzer (+)       ; (-) -> riel (-)
   SDA -------> OLED SDA
   SCL -------> OLED SCL
   |
   HC-SR04 VCC -> riel (+)   ;  GND -> riel (-)
   OLED    VCC -> riel (+)   ;  GND -> riel (-)

   [R] = resistencia 220-330 ohm en serie con cada LED
\end{lstlisting}

% ===================== NOTAS =====================
\section{Notas de armado}
\begin{itemize}
\item \textbf{Polaridad del LED:} la \textbf{pata larga} es el ánodo (+) y va hacia la
resistencia/pin; la \textbf{pata corta} es el cátodo ($-$) y va al riel azul. Si está al revés,
no enciende (no se daña).
\item \textbf{Resistencia en serie:} cada LED lleva su resistencia (220--330\,$\Omega$). Nunca
conectes un LED directo a un pin sin resistencia.
\item \textbf{Buzzer:} tiene polaridad; la pata larga (o el lado marcado \texttt{+}) va a
\pin{D8}. Si es del tipo sellado, retira el sticker que tapa el agujero del sonido.
\end{itemize}

\begin{ojo}
\textbf{Echo del HC-SR04 va directo a \pin{D10}.} La UNO (R3/R4) trabaja a 5\,V, así que no
necesita divisor de voltaje. (Eso solo haría falta en una placa de 3.3\,V.)
\end{ojo}

\begin{ojo}
\textbf{OLED en negro pero el resto funciona:} casi siempre es \textbf{SDA y SCL cruzados} o que
la dirección I2C es \texttt{0x3D} en vez de \texttt{0x3C}. Corrige el cruce y reinicia la placa,
o cambia \texttt{DIR\_OLED} en el firmware.
\end{ojo}

% ===================== FIRMWARE =====================
\section{Firmware}
El programa está en \texttt{arduino-r4/geogreen\_proto/} (proyecto PlatformIO). Para grabarlo en
la placa, con ella conectada por USB:

\begin{lstlisting}[style=cmd]
pio run -d arduino-r4/geogreen_proto -t upload
\end{lstlisting}

Al arrancar hace un \textbf{autotest}: enciende verde $\rightarrow$ amarillo $\rightarrow$ rojo,
uno por uno (confirma el cableado de los LEDs). La matriz pasa la marca \textbf{GeoGreen · AIEP}
en bucle.

\subsection{Perillas de calibración (constantes arriba del archivo)}
\begin{tabularx}{\linewidth}{@{}L{4.2cm}L{2.2cm}Y@{}}
\hline
\textbf{Constante} & \textbf{Valor} & \textbf{Qué ajusta}\\
\hline
\texttt{DIST\_VACIO} & 25 cm & distancia con el contenedor vacío (0\%)\\
\texttt{DIST\_LLENO} & 3 cm & distancia con el contenedor lleno (100\%)\\
\texttt{UMBRAL\_MEDIO} & 40\% & corte verde $\rightarrow$ amarillo\\
\texttt{UMBRAL\_ALTO} & 80\% & corte amarillo $\rightarrow$ rojo + buzzer\\
\texttt{N\_MUESTRAS} & 7 & lecturas por ciclo (mediana)\\
\texttt{SALTO\_MAX} / \texttt{CONFIRMAR} & 6 cm / 3 & rechazo de saltos falsos del sensor\\
\texttt{DEADBAND\_PCT} & 1.5 & cuánto debe cambiar para refrescar el \%\\
\hline
\end{tabularx}

\begin{nota}
\textbf{Para un contenedor real} (no demo de escritorio), sube \texttt{DIST\_VACIO} a la
profundidad real del contenedor (p.ej. 100 cm) y deja \texttt{DIST\_LLENO} en la distancia mínima
útil del sensor.
\end{nota}

% ===================== PUESTA EN MARCHA =====================
\section{Puesta en marcha}
\begin{enumerate}
\item Arma el cableado según el \textbf{mapa de pines}.
\item Conecta la placa por USB y graba el firmware.
\item \textbf{Verifica el autotest:} deben encender los 3 LEDs en orden. Si uno no enciende, está
al revés.
\item Mueve la mano frente al sensor: el semáforo y la OLED deben cambiar
verde $\rightarrow$ amarillo $\rightarrow$ rojo, y el buzzer sonar en rojo.
\item La matriz debe estar pasando \textbf{GeoGreen · AIEP}.
\end{enumerate}

\end{document}
